\documentclass[11pt]{article}
\usepackage[letterpaper,margin=1.05in]{geometry}
\usepackage[T1]{fontenc}
\usepackage{lmodern,amsmath,amssymb,amsthm,microtype}
\usepackage[round,authoryear]{natbib}
\usepackage[colorlinks=true,linkcolor=blue,citecolor=blue,urlcolor=blue]{hyperref}
\newtheorem{theorem}{Theorem}
\newtheorem{proposition}[theorem]{Proposition}
\newtheorem{lemma}[theorem]{Lemma}
\newtheorem{corollary}[theorem]{Corollary}
\theoremstyle{definition}
\newtheorem{definition}[theorem]{Definition}
\newtheorem{remark}[theorem]{Remark}
\newtheorem{example}[theorem]{Example}
\newcommand{\GH}{\mathsf{GH}}
\newcommand{\Mink}{\mathbb M}
\newcommand{\U}{\mathcal U}
\newcommand{\Ifut}{I^{+}}
\newcommand{\Ichr}{I^{-}}
\newcommand{\Jp}{J^{+}}
\newcommand{\Jm}{J^{-}}
\newcommand{\B}{\mathcal B}
\title{Stein's Theorem in General Relativity:\\[2pt]
\large locality, not symmetry, fixes what has become}
\author{Working manuscript}
\date{24 October 2026}
\begin{document}
\maketitle

\begin{abstract}
Stein proved that in Minkowski spacetime a transitive relation of becoming
that is invariant under the spacetime's symmetries, and that makes some past
event definite for an event, must be the causal past or degenerate. Clifton
and Hogarth showed that, assuming only that the chronological past has become,
the options are the causal past, the chronological past and the universal
relation; they extended the result to worldline-dependent becoming and
remarked that there is ``no
obvious way'' to extend it beyond Minkowski spacetime. We give such an
extension and locate exactly what it depends on. In a single curved
spacetime, symmetry-invariance is too weak: generic spacetimes have no
symmetries, and transitive relations containing the chronological past can
then relate spacelike events. The right replacement is local covariance in
the sense of locally covariant quantum field theory: one rule of becoming for
all globally hyperbolic spacetimes, compatible with isometric embeddings of
causally convex regions. Under that hypothesis we prove, by a gluing
argument, that either everything has become for everything in every
spacetime, or in every spacetime becoming lies between the chronological and
the causal past; the worldline-dependent form follows. The residual freedom is exactly the status of
lightlike-separated pairs, and we classify it: it can depend on curvature
along the connecting light ray, so uniqueness fails even for relations
determined by causal structure alone; topological regularity removes it. In
causal set theory the same hypothesis yields an exact Stein theorem,
although the analogue fails, as Wüthrich and Callender showed, under
invariance within a single causal set. Locality also isolates
Gödel-type universes, so that universal becoming there constrains nothing
here. Finally we show why the result does not transfer routinely to
spacetimes restricted by energy conditions, and isolate the open problem on
which such a transfer depends.
\end{abstract}

\section{Introduction}\label{sec:intro}

Stein's theorem is the strongest positive result in the debate over
relativistic becoming \citep{Stein1968,Stein1991}. Read as a constraint, it
says that an objective, transitive becoming relation adapted to Minkowski
geometry can only be past causal connectibility, apart from trivial or
universal relations. Clifton and Hogarth gave the result its now-standard form
\citep{CH1995}: their Theorem~1 assumes invariance under the
time-orientation-preserving isometries of Minkowski spacetime, transitivity,
and \emph{chronological becoming} (every event in an event's chronological
past has become for it), and concludes that becoming is $\Ichr$, $\Jm$ or the
universal relation $\U$. Their Theorem~2 extends this to becoming relative to
an observer's worldline. They end with a problem:
\begin{quote}
each general relativistic spacetime has its own metric-preserving
automorphisms and each its own distinctive geometry, leaving no obvious way
to extend Theorem 1 beyond Minkowski spacetime, much less to extend Theorem 2
\citep[p.~385]{CH1995}.
\end{quote}
They call the question of becoming ``in the real world'' ``more difficult,
though more interesting'' (ibid., p.~386).

The difficulty is real. Stein's argument runs entirely on symmetry: once one
spacelike pair is related, boosts and rotations relate them all, and
transitivity collapses the relation to $\U$. A generic curved spacetime has
no isometries at all, so the corresponding single-spacetime statement is
simply false (Proposition~\ref{prop:single}). Wüthrich and Callender observed
the analogous failure in causal set theory, where invariance under the
automorphisms of one causal set permits exotic becoming relations
\citep{WC2017}.

Our diagnosis is that the analogue of Minkowskian symmetry in general
relativity is not invariance within a spacetime but \emph{uniformity and
locality across spacetimes}: one rule, applied in every globally hyperbolic
spacetime, whose verdict on a pair of events depends only on a causally
convex region containing them. This is the hypothesis of local covariance
used to say when a quantum field theory is ``the same in all spacetimes''
\citep{BFV2003,FV2012}, and it is closely related to the locality
intuition that Viglione defends informally against cosmic simultaneity
\citep{Viglione2026}. With it, Stein's theorem has a curved-spacetime form:

\begin{itemize}
\item \textbf{No becoming outside the causal past} (Theorem~\ref{thm:main}).
Either becoming is universal in every spacetime, or in every spacetime it lies
between $\Ichr$ and $\Jm$. The worldline-dependent version follows
(Corollary~\ref{cor:worldline}).
\item \textbf{The residual freedom is the light cone} (Theorem~\ref{thm:class},
Theorem~\ref{thm:weyl}, Proposition~\ref{prop:reg}). What remains open is
exactly which lightlike-separated pairs count. We classify the possibilities;
they can depend on Weyl curvature along the connecting light ray, so
uniqueness fails even for relations fixed by causal structure. Closedness or
openness of the relation selects $\Jm$ or $\Ichr$.
\item \textbf{Locality is indispensable} (Propositions~\ref{prop:single} and
\ref{prop:nonlocal}). Natural but non-local rules, such as ordering by a
canonical cosmological time, satisfy every other condition and make spacelike
events become.
\item \textbf{Causal sets} (Proposition~\ref{prop:causets}). Under the same
hypothesis Stein's theorem holds exactly: becoming is the causal order or
universal. Discreteness removes the light-cone freedom.
\item \textbf{Gödel universes are isolated} (Proposition~\ref{prop:godel}).
Totally vicious spacetimes have no proper causally convex regions, so local
covariance transmits nothing from them.
\item \textbf{Energy conditions} (Section~\ref{sec:energy}). The proof uses
arbitrary globally hyperbolic metrics. Under the dominant energy condition the
key construction is provably unavailable, and the transfer reduces to an open
gluing problem that we state.
\end{itemize}

The philosophical upshot is stated in Section~\ref{sec:phil}. What Stein's
theorem tells us about relativity survives in general relativity in a robust
and a fragile part. The robust part is that nothing outside an event's causal
past can have become for it unless everything has. The fragile part is the
uniqueness of the answer, which was an artefact of Minkowskian symmetry.

\section{Framework}\label{sec:framework}

Spacetimes are smooth, connected, four-dimensional, oriented and
time-oriented Lorentzian manifolds. We follow Clifton and Hogarth in taking
$\Ichr,\Jm$ to be reflexive \citep[p.~360]{CH1995}, write $x\ll y$ for
$x\in\Ichr(y)$, $x\neq y$, and $x\le y$ for $x\in\Jm(y)$. We write
$\Ifut_{\mathrm s},\Ichr_{\mathrm s}$ for the strict chronological future and
past (defined by timelike curves, so $x\in\Ichr_{\mathrm s}(y)\iff x\ll y$ in
causal spacetimes); thus $\Ichr=\Ichr_{\mathrm s}\cup\Delta$. Standard causal facts
are cited from Minguzzi and Sánchez \citep{MS2008}, abbreviated MS, whose
$I^\pm$ are our strict sets.

\begin{definition}
$\GH$ is the category whose objects are globally hyperbolic spacetimes and
whose morphisms $\psi:M\to M'$ are smooth isometric embeddings preserving
orientation and time orientation, with $\psi(M)$ open and causally convex in
$M'$. This is the category $\mathsf{Loc}$ of locally covariant quantum field
theory \citep[\S2]{BFV2003}, \citep[\S2.2]{FV2012}.
\end{definition}

\begin{definition}
A \emph{becoming theory} assigns to each $M\in\GH$ a relation
$B_M\subseteq M\times M$ (read $x\,B_M\,b$ as ``$x$ has become for $b$'') such
that, for every $M$:
\begin{enumerate}
\item[(LC)] for every morphism $\psi:M\to M'$, $x\,B_M\,b\iff
\psi(x)\,B_{M'}\,\psi(b)$;
\item[(T)] $B_M$ is transitive;
\item[(CB)] $\Ichr_M\subseteq B_M$.
\end{enumerate}
\end{definition}

(CB) is Clifton and Hogarth's chronological becoming; they derive it from
``worldline becoming'' and ``worldline transitivity'' \citep[Lemma,
p.~364]{CH1995}, and the derivation carries over verbatim to curved
spacetime. (LC) contains invariance under every orientation- and
time-orientation-preserving isometry of each $M$, which on Minkowski spacetime
$\Mink$ is the proper orthochronous Poincaré group. Its further content is
\emph{locality}: if $N\subseteq M$ is open, causally convex and contains $x$
and $b$, then whether $x$ has become for $b$ is settled by $N$ alone.

\begin{remark}[The off-shell premise]\label{rem:offshell}
$\GH$ contains every globally hyperbolic metric, not only solutions of
Einstein's equations; the category is defined kinematically in
\citep{BFV2003,FV2012}. Theorem~\ref{thm:main} uses this: its proof builds
auxiliary spacetimes that need not satisfy any field equation with reasonable
matter. This is natural for a kinematic notion of becoming, defined (as Stein
and Clifton--Hogarth define theirs) from geometry alone, but it is a premise.
Section~\ref{sec:energy} examines what happens without it.
\end{remark}

\section{Why one spacetime is not enough}\label{sec:single}

\begin{proposition}\label{prop:single}
Let $M$ be globally hyperbolic and $x,b\in M$ spacelike-separated. Then
$R=\Jm\cup(\Jm(x)\times\Jp(b))$, where $(e,z)$ in the second set means $e$
has become for $z$, is reflexive and transitive, contains $\Ichr$, relates $x$
to $b$, and is neither $\Ichr$, $\Jm$ nor $\U$. If $M$ has no nontrivial
isometries, $R$ is invariant under all of them.
\end{proposition}

\begin{proof}
Write $P=\Jm(x)\times\Jp(b)$. Compositions $\Jm\circ\Jm$, $\Jm\circ P$ and
$P\circ\Jm$ land in $\Jm$ or $P$ because $\le$ is transitive. A composition
$P\circ P$ would need some $q$ with $b\le q\le x$, which is impossible. The
pair $(b,x)$ is not in $R$, so $R\neq\U$.
\end{proof}

So for a single generic spacetime, Stein-type uniqueness fails completely; the
dust-cosmology ordering by scalar curvature (in $a(t)=t^{2/3}$ dust, curvature
fixes cosmic time) is a symmetric instance. The next example shows that
uniformity across spacetimes is still not enough without locality.

\begin{proposition}\label{prop:nonlocal}
Let $\tau_M(p)$ be the supremum of the lengths of past-directed causal curves
from $p$ (the cosmological time of Andersson, Galloway and Howard
\citep{AGH1998}). Put $B_M=\{(x,b):\tau_M(x)\le\tau_M(b)\}$ if $\tau_M$ is
finite on $M$, and $B_M=\Jm_M$ otherwise. Then $B$ is invariant under all
isometries and under the morphisms of $\GH$ that are onto, each $B_M$ is
reflexive and transitive and satisfies (CB), and $B$ relates spacelike pairs
in every spacetime where $\tau$ is finite (e.g.\ big-bang
Friedmann--Lemaître--Robertson--Walker models). It is not a becoming theory:
it violates (LC).
\end{proposition}

\begin{proof}
Isometries preserve lengths of causal curves. If $x\ll b$, concatenating
past curves from $x$ with a timelike curve from $x$ to $b$ gives
$\tau(b)\ge\tau(x)+d(x,b)>\tau(x)$, so (CB) holds; transitivity is that of
$\le$ on the reals. Failure of (LC): the Milne wedge $W=\Ifut_\Mink(0)$ is open,
causally convex and globally hyperbolic, and its inclusion into $\Mink$ is a
morphism. On $W$, $\tau_W(p)$ is the proper time from the origin, finite, so
$B_W$ relates distinct points of the same hyperboloid $\{\tau_W=c\}$, which are
spacelike-separated; on $\Mink$, $\tau=\infty$ and $B_\Mink=\Jm$ does not relate
them.
\end{proof}

\section{Results}\label{sec:results}

\begin{theorem}[Clifton--Hogarth]\label{thm:flat}
For every becoming theory, $B_\Mink\in\{\Ichr,\Jm,\U\}$.
\end{theorem}

\begin{proof}
Proper orthochronous Poincaré maps are automorphisms of $\Mink$ in $\GH$, so
$B_\Mink$ is invariant under them by (LC). Clifton and Hogarth's Theorem~1
\citep[p.~365]{CH1995} gives the conclusion from this invariance,
transitivity and chronological becoming. Its proof (pp.~365--367) uses only
point-preserving spatial rotations, boosts and translations, and they note
that it does not use spatial reflections (pp.~385--386). Their stated
hypothesis is invariance under all time-orientation-preserving isometries,
reflections included; using only the proper orthochronous subgroup is a
weakening whose adequacy rests on their proof, not their statement.
\end{proof}

\begin{theorem}[No becoming outside the causal past]\label{thm:main}
Let $B$ be a becoming theory.
\begin{enumerate}
\item[(a)] If $B_\Mink\neq\U$, then $B_M\subseteq\Jm_M$ for every $M\in\GH$.
\item[(b)] If $B_\Mink=\U$, then $B_M=\U_M$ for every $M\in\GH$.
\end{enumerate}
Hence either $B$ is universal everywhere, or $\Ichr_M\subseteq
B_M\subseteq\Jm_M$ for every $M$.
\end{theorem}

\begin{corollary}[Worldline-dependent becoming]\label{cor:worldline}
Suppose that for each $M\in\GH$, each inextendible timelike curve $\lambda$ in
$M$ and each $b\in\lambda$ a set $B_{M,\lambda}(b)$ is given, satisfying
Clifton--Hogarth's worldline becoming ($\Ichr(b)\cap\lambda\subseteq
B_{M,\lambda}(b)$) and worldline transitivity ($a\in B_{M,\lambda}(b)$ and
$b\in B_{M,\lambda'}(c)$ imply $a\in B_{M,\lambda'}(c)$) \citep[pp.~362--363]{CH1995},
and local covariance in the form: for a morphism $\psi:N\to M$, an inextendible
timelike $\lambda_N$ in $N$ and any inextendible timelike $\lambda$ in $M$
extending $\psi\circ\lambda_N$, $a\in B_{N,\lambda_N}(b)\iff\psi(a)\in
B_{M,\lambda}(\psi(b))$. Then $B_{M,\lambda}(b)$ does not depend on $\lambda$, and
the resulting relation is a becoming theory. Hence Theorem~\ref{thm:main}
applies to worldline-dependent becoming.
\end{corollary}

Clifton and Hogarth anticipated that extending their Theorem~2 would be
harder still (``much less to extend Theorem 2'', p.~385). Under local
covariance it is not: as they themselves note (p.~363), worldline transitivity
makes observers whose worldlines cross agree, and the remaining work is done by
Theorem~\ref{thm:main}.

Write $E^-=\Jm\setminus\Ichr$ for the horismos (lightlike-only relation).

\begin{theorem}[Classification of the residual freedom]\label{thm:class}
Let $B$ be a non-universal becoming theory and $N_M=B_M\cap E^-_M$. Then $N$
is locally covariant and satisfies
\begin{enumerate}
\item[(NC)] if $(x,y),(y,b)\in N_M$ and $(x,b)\in E^-_M$, then $(x,b)\in N_M$.
\end{enumerate}
Conversely, if $N\subseteq E^-$ is locally covariant and satisfies (NC), then
$B=\Ichr\cup N$ is a becoming theory. On $\Mink$, $N$ is empty or all of $E^-$.
\end{theorem}

\begin{theorem}[Non-uniqueness]\label{thm:weyl}
Let $N_M(x,b)$ hold iff $x\in E^-(b)$, $x\neq b$, and some lightlike geodesic
segment from $x$ to $b$ contains a point at which the Weyl tensor is nonzero.
Then $\Ichr\cup N$ is a becoming theory equal to $\Ichr$ on $\Mink$, equal to
$\Jm$ on any spacetime with nowhere-vanishing Weyl tensor (such as the
Schwarzschild exterior), and strictly between $\Ichr$ and $\Jm$ on spacetimes
that are flat on a nonempty causally convex open set and have nonzero Weyl
tensor somewhere. It depends only on the conformal structure (not merely the
curvature: conformally flat spacetimes such as FLRW give $\Ichr$).
\end{theorem}

\begin{proposition}[Regularity]\label{prop:reg}
Let $B$ be non-universal. If every $B_M$ is closed in $M\times M$, then
$B=\Jm$. If every $B_M\setminus\Delta_M$ is open, then $B=\Ichr$.
\end{proposition}

\section{Proofs}\label{sec:proofs}

\subsection{Theorem \ref{thm:class}, Theorem \ref{thm:weyl}, Proposition \ref{prop:reg}}

We use: $\Ifut_{\mathrm s}$ is open in $M\times M$ (MS Prop.~2.16);
$\overline{\Jp}=\overline{\Ifut_{\mathrm s}}$ and $\operatorname{Int}\Jp=
\Ifut_{\mathrm s}$ (MS Prop.~2.17), and likewise for the pasts; if $p\ll r\le q$
or $p\le r\ll q$ then $p\ll q$, and if $r\in E^+(p)$, $q\in E^+(r)$ then
$q\in E^+(p)$ or $p\ll q$ (MS Prop.~2.31); a causal curve that is not an
unbroken lightlike pregeodesic can be deformed to a timelike curve with the
same endpoints (MS Thm.~2.30).

\begin{proof}[Proof of Theorem \ref{thm:class}]
By Theorem~\ref{thm:main}(a), $B=\Ichr\cup N$. Timelike and causal curves
between points of a causally convex set stay in it, so $\Ichr$, $\Jm$ and $E^-$
are locally covariant and $N$ inherits (LC); (NC) is transitivity restricted
to horismotic pairs. Conversely, $\Ichr\cup N$ is reflexive and satisfies (CB);
for transitivity, a composition with a chronological step is chronological by
MS 2.31, and a composition of two $N$-steps is chronological or horismotic, in
which case (NC) applies. The last claim is Theorem~\ref{thm:flat}.
\end{proof}

\begin{proof}[Proof of Theorem \ref{thm:weyl}]
(NC): if $x\to y\to b$ along lightlike segments $\gamma_1,\gamma_2$ and
$x\in E^-(b)$, then $\gamma_1\gamma_2$ is a causal curve between horismotic
points, hence an unbroken lightlike pregeodesic (MS Thm.~2.30), and it
contains the Weyl point of $\gamma_1$. (LC): lightlike geodesics between points
of a causally convex image stay in it, and isometries preserve the Weyl
tensor. On $\Mink$ the Weyl tensor vanishes. If it vanishes nowhere, every
horismotic pair is in $N$. If $M$ is flat on an open causally convex $V$ and
the Weyl tensor is nonzero at some point $z$, horismotic pairs inside $V$ are
not in $N$ (every causal curve between them stays in $V$), while a horismotic
pair joined by a lightlike segment through a neighbourhood of $z$ is. The Weyl
tensor, as a $(1,3)$ tensor, is invariant under conformal rescaling (standard),
and unparametrized lightlike pregeodesics are conformally invariant (MS
\S2.6).
\end{proof}

\begin{proof}[Proof of Proposition \ref{prop:reg}]
If $B_M$ is closed and contains $\Ichr$, it contains
$\overline{\Ichr}=\overline{\Ichr_{\mathrm s}}=\overline{\Jm}\supseteq\Jm$, so
$B=\Jm$ by Theorem~\ref{thm:main}(a). If $B_M\setminus\Delta$ is open and
contained in $\Jm$, it lies in $\operatorname{Int}\Jm=\Ichr_{\mathrm s}$, so
$B_M\subseteq\Ichr_{\mathrm s}\cup\Delta=\Ichr$, and (CB) gives $B=\Ichr$.
\end{proof}

\begin{proof}[Proof of Corollary \ref{cor:worldline}]
By worldline becoming, $b\in B_{M,\lambda'}(b)$ for every $\lambda'\ni b$; with
$c=b$, worldline transitivity gives $B_{M,\lambda}(b)\subseteq B_{M,\lambda'}(b)$,
and by symmetry they are equal. Call the common set $B_M(b)$. Transitivity: if
$a\in B_M(b)$ and $b\in B_M(c)$, choose $\lambda\ni b$, $\lambda'\ni c$ and apply
worldline transitivity. (CB): if $a\ll b$, a timelike curve from $a$ to $b$
extends to an inextendible timelike curve $\lambda$, and $a\in\Ichr(b)\cap\lambda
\subseteq B_M(b)$. (LC): for $\psi:N\to M$ take $\lambda_N\ni b$ and an
inextendible extension $\lambda$ of $\psi\circ\lambda_N$; then $a\in B_N(b)\iff
a\in B_{N,\lambda_N}(b)\iff\psi(a)\in B_{M,\lambda}(\psi(b))\iff\psi(a)\in
B_M(\psi(b))$.
\end{proof}

\subsection{Theorem \ref{thm:main}}

The idea: a forbidden pair in a curved spacetime is moved, by local
covariance, into an auxiliary spacetime that agrees with the original near
the pair and is flat elsewhere; chronological becoming and transitivity then
produce a forbidden pair in a flat region, contradicting
Theorem~\ref{thm:flat}.

\begin{lemma}[Future pairs reduce to spacelike pairs]\label{lem:R}
If $x\in\Jp(b)\setminus\{b\}$, there is $e\ll x$ spacelike to $b$.
\end{lemma}

\begin{proof}
Let $U$ be a convex normal neighbourhood of $b$ and $V\subseteq U$ an open
neighbourhood of $b$ that is causally convex in $M$ ($M$ is strongly causal;
MS Remark~2.12 and the construction following it). Causal relations in $U$,
regarded as a spacetime, are read off $\exp^{-1}$ (MS Prop.~2.10), and for
points of $V$ they agree with those of $M$ (MS Remark~2.12). A future causal
curve from $b$ to $x$ is nonconstant ($M$ is causal), so it contains some
$b_1\in V\setminus\{b\}$; then $b\le b_1\le x$ and $u=\exp_b^{-1}(b_1)$ is
a nonzero future causal vector. Choose a spacelike $a\in T_bM$ with
$g(u,a)>0$ and put $q_\varepsilon=\exp_b(\varepsilon a)$. Synge's world
function $\sigma$ is smooth on $U\times U$ with $\nabla_q\sigma(q,b_1)|_{q=b}=-u$,
so $2\sigma(q_\varepsilon,b_1)=g(u,u)-2\varepsilon g(u,a)+O(\varepsilon^2)<0$ for
small $\varepsilon>0$. The vector $\exp^{-1}_{q_\varepsilon}(b_1)$ depends
continuously on $\varepsilon$ and tends to $u\neq0$, so for small $\varepsilon$
it is future-directed timelike: $q_\varepsilon\ll b_1$. Also
$2\sigma(q_\varepsilon,b)=\varepsilon^2g(a,a)>0$, so $q_\varepsilon$ is spacelike
to $b$. By MS Prop.~2.31(i), $e=q_\varepsilon\ll x$.
\end{proof}

The same local argument shows that every point has spacelike-separated
points arbitrarily close to it.

\begin{lemma}[Spacelike chains]\label{lem:E}
In connected globally hyperbolic $M$, any two points are joined by a finite
chain of consecutively spacelike-separated points.
\end{lemma}

\begin{proof}
Globally hyperbolic spacetimes (MS Def.~3.70) are causally simple (MS
Prop.~3.71), so $\Jp(z)$ and $\Jm(z)$ are closed and the set of points
spacelike to $z$ is open. Fix $p$ and let $S$ be the set of points reachable
from $p$. $S$ is open: if $q\in S$, pick $z$ spacelike to $q$; a
neighbourhood of $q$ is spacelike to $z$ and reachable via $q,z$.
$M\setminus S$ is open: if $q\notin S$ and $c$ is spacelike to $q$, a
neighbourhood $W$ of $q$ is spacelike to $c$; were some $q''\in W$ in $S$, then
$c\in S$ and $q\in S$. $M$ is connected.
\end{proof}

The example $M=\mathbb R\times S^3$, $b=(0,p)$, $x=(3\pi,p)$ shows that a single
point spacelike to two causally related points need not exist; the chain is
needed.

\paragraph{Set-up.} Fix spacelike-separated $x,b\in M$. Since $J^+\subseteq
M\times M$ is closed (MS Prop.~3.68(iii), 3.71), there are small closed
spacelike embedded $3$-discs $D_x\ni x$, $D_b\ni b$, each acausal, with no causal
curve between them; so $D_x\cup D_b$ is a compact acausal spacelike
submanifold with boundary, and Bernal and Sánchez give a smooth spacelike
Cauchy hypersurface $S\supseteq D_x\cup D_b$ \citep[Thm.~1.1, Rem.~4.14]{BS2006}.
Choose $\delta_0>0$ and a smooth embedded curve $c:[-\delta_0,L+\delta_0]\to S$,
unit speed for the induced metric $h$, with $c(0)=x$, $c(L)=b$.

\begin{lemma}[Adapted coordinates]\label{lem:B}
Let $n$ be the future unit normal to $S$, $(e_1=c',e_2,e_3)$ an $h$-orthonormal
frame along $c$ with $(n,e_1,e_2,e_3)$ positively oriented,
$\sigma(s,y)=\exp^h_{c(s)}(y_2e_2+y_3e_3)$, and
$\Phi(t,s,y)=\exp^g_{\sigma(s,y)}(t\,n)$. Write
$T_\rho=\{|t|<\rho,\ -\rho<s<L+\rho,\ |y|<\rho\}$. For some $\rho_0\le\delta_0$,
$\Phi$ is a diffeomorphism from $\overline{T_{\rho_0}}$ onto a closed
neighbourhood of $c([0,L])$, orientation- and time-orientation-preserving;
$\Phi(\{t=0\})\subseteq S$ while $\Phi(t,\cdot)\in\Ifut(S)$ for $t>0$ and
$\in\Ichr(S)$ for $t<0$; $\Phi^*g=-dt^2+h_t(s,y)$; and
$\|\Phi^*g-\eta\|\le K(|t|+|y|)$ on $\overline{T_{\rho_0}}$.
\end{lemma}

\begin{proof}
Gaussian normal coordinates to $S$ composed with Riemannian normal
coordinates of $S$ along the compact embedded arc; points $\Phi(t,\cdot)$ lie
on the normal timelike geodesic through a point of the acausal $S$. At
$(t,y)=(0,0)$ the metric is exactly $\eta$ because the frame is orthonormal on
$c$; the estimate is smoothness on a compact set.
\end{proof}

\begin{lemma}[Cone control]\label{lem:C}
Let $\eta_c=-c^2dt^2+ds^2+|dy|^2$ and $\tfrac12<\mu<1<\lambda$. If
$\|h-\eta_1\|\le\varepsilon<\min\{(1-\mu^2)/(1+\mu^2),(\lambda^2-1)/(\lambda^2+1)\}$
(operator norm), then $h$ is Lorentzian with $\partial_t$ timelike, and for
$v=(1,u)$: $|u|\le\mu\Rightarrow h(v,v)<0$, $|u|\ge\lambda\Rightarrow
h(v,v)>0$. The bound is preserved by $\chi h+(1-\chi)\eta_1$, $0\le\chi\le1$.
\end{lemma}

\begin{proof}
$h(v,v)\le-1+|u|^2+\varepsilon(1+|u|^2)$ and $h(v,v)\ge(1-\varepsilon)|u|^2-
(1+\varepsilon)$. The spatial block is $\ge(1-\varepsilon)I>0$ and
$h(\partial_t,\partial_t)<0$, which forces signature $(-,+,+,+)$.
\end{proof}

\paragraph{Constants and the auxiliary spacetime.} In this order: fix
$\mu,\lambda,\varepsilon$ as in Lemma~\ref{lem:C}; Lemma~\ref{lem:B} gives
$\rho_0$ and $K$; choose $r>0$ with $3r<\rho_0$, $6Kr\le\varepsilon$ and
$8\lambda r<L$. Then $\overline{T_{2r}}\subset T_{3r}$; let $\chi$ be smooth,
$\chi=1$ on $\overline{T_{2r}}$, $\operatorname{supp}\chi\subset T_{3r}$, and on
$\mathbb R^4$ put $g'=\chi\,\Phi^*g+(1-\chi)\eta$. On $T_{3r}$,
$\|\Phi^*g-\eta\|\le6Kr\le\varepsilon$, so Lemma~\ref{lem:C} applies to $g'$
everywhere and to $\Phi^*g$ on $T_{3r}$. Every $g'$-causal vector therefore
has $|d(s,y)/dt|\le\lambda$; $t$ increases strictly along future causal
curves; and an inextendible causal curve along which $t$ had a finite limit
would have convergent spatial coordinates and extend. Hence every level
$\{t=c\}$ is a Cauchy hypersurface and $M'=(\mathbb R^4,g')$ is globally
hyperbolic (MS Thm.~3.75), flat outside $T_{3r}$, and
$\psi=\Phi^{-1}:\Phi(T_{2r})\to T_{2r}$ is an isometry preserving orientation
and time orientation.

\begin{lemma}[The slab region is self-determined]\label{lem:D}
Let $B_0=\{s\in[-r/2,L+r/2],\ |y|\le r/2\}$, $\Sigma=\Phi(\{0\}\times B_0)\subset
S$ and $\Sigma'=\{0\}\times B_0$. Then:
\begin{enumerate}
\item[(i)] $D_M(\Sigma)\subseteq\Phi(T_{2r})$ and $D_{M'}(\Sigma')\subseteq T_{2r}$;
\item[(ii)] $\psi(D_M(\Sigma))=D_{M'}(\Sigma')$;
\item[(iii)] $N=\operatorname{int}D_M(\Sigma)$ is causally convex in $M$,
globally hyperbolic, connected, and contains $x,b$; likewise
$\psi(N)=\operatorname{int}D_{M'}(\Sigma')$ in $M'$.
\end{enumerate}
\end{lemma}

\begin{proof}
\emph{Escape paths.} Let $q\in\partial T_{2r}$ with $0<t(q)\le2r$. If the
projection $(s_q,y_q)$ is not in $B_0$, follow the vertical segment down to
$t=0$. Otherwise $q$ lies on the lid $t=2r$ with $|y_q|\le r/2$; move with
velocity $(dt,dy)=(-1,\tfrac\mu2\hat y)$ for coordinate time $2r$, where $\hat y$
is a unit vector ($y_q/|y_q|$ if $y_q\neq0$), ending at $|y|\in(r/2,3r)$
because $\mu r>r/2$ and $r/2+\mu r<3r$. Either path lies in $T_{3r}$, has
coordinate speed $<\mu$, so is timelike for $\Phi^*g$ and for $g'$
(Lemma~\ref{lem:C}), and ends at a point of $\{t=0\}$ outside $\Sigma'$, i.e.\ on
$S\setminus\Sigma$ in $M$ ($\Phi$ is injective). A past-directed causal curve
starting on a Cauchy hypersurface never returns to it. The time-reversed paths
serve points with $t(q)<0$.

(i) Let $p\in D^+_M(\Sigma)$ with $p\notin\Phi(T_{2r})$, and let $\gamma$ be a
past-inextendible causal curve from $p$. It meets $\Sigma$; let $\sigma_1$ be
its first point on $\Sigma$, and $q$ the last point of $\partial\Phi(T_{2r})$ on
the segment from $p$ to $\sigma_1$. After $q$ the segment lies in
$\Phi(T_{2r})$, where $t$ increases along future causal curves, so $t(q)>0$.
The segment $p\to q$, the escape path from $q$, and any past-inextendible
continuation form a past-inextendible causal curve from $p$ missing $\Sigma$, a
contradiction. Time-reverse for $D^-$. The same argument works in $M'$, where
$\{t=0\}$ is a Cauchy hypersurface.

(ii) Let $p\in D^+_M(\Sigma)$, so $p\in\Phi(T_{2r})$ by (i), and let
$\gamma'$ be a past-inextendible $g'$-causal curve from $\psi(p)$. If
$\psi(p)\in\Sigma'$ there is nothing to prove; otherwise $t(\psi(p))>0$. Along
$\gamma'$, $t\to-\infty$; let $\tau_0$ be the first parameter at which $t=0$ and
$\tau_{\rm exit}$ the first at which $\gamma'$ meets $\partial T_{2r}$. The chart
and the agreement $g'=\Phi^*g$ extend to a neighbourhood of
$\overline{T_{2r}}$, so segments of $\gamma'$ up to either time pull back by
$\Phi$. If $\tau_0\le\tau_{\rm exit}$ (equality included): were
$\gamma'(\tau_0)\notin\Sigma'$, the pulled-back segment ends on $S\setminus\Sigma$
and, continued past $S$, is a past-inextendible $M$-curve from $p$ missing
$\Sigma$ (points with $t>0$ are not on $S$). If $\tau_{\rm exit}<\tau_0$, then
$t>0$ at the exit point, and the pulled-back segment followed by the escape
path in $M$ again misses $\Sigma$. Both contradict $p\in D^+_M(\Sigma)$, so $\gamma'$
meets $\Sigma'$. The reverse inclusion is the same argument with $M$ and $M'$
exchanged, using $D_{M'}(\Sigma')\subseteq T_{2r}$: no global coordinate $t$ on
$M$ is needed, since a past-inextendible $M$-curve from a point of
$\Phi(T_{2r})$ either leaves $\Phi(T_{2r})$ or reaches the Cauchy surface $S$,
and the same two cases apply. $D^-$ is time-reversed.

(iii) For $q\in D^+(\Sigma)$, $\Jm(q)\cap\Jp(\Sigma)\subseteq D^+(\Sigma)$: a
past-inextendible curve from $r$ in this set, prefixed by a causal curve
$q\to r$, meets $\Sigma$; if only on the prefix, at $z$, then $\sigma\le r\le z$
for some $\sigma\in\Sigma$, and acausality gives $r=z\in\Sigma$. Hence $D(\Sigma)$
is causally convex: if $p,q\in D(\Sigma)$, $r$ lies on a causal curve $p\to q$,
and an inextendible causal curve $\lambda$ through $r$ misses $\Sigma$, then in
the mixed case $p\in D^-$, $q\in D^+$ the curves $(p\to r)\cdot\lambda^+$ and
$\lambda^-\cdot(r\to q)$ meet $\Sigma$ on $p\to r$ and on $r\to q$, and
acausality forces both meeting points to be $r$; the cases $p,q\in D^\pm$
follow from the inclusion just proved. For interior points choose $p^-\ll p$
and $q^+\gg q$ in $\operatorname{int}D$; then
$J(p,q)\subseteq\Ifut_{\mathrm s}(p^-)\cap\Ichr_{\mathrm s}(q^+)\subseteq
D(\Sigma)$, an open set (MS Props.~2.16, 2.31), so $N$ is causally convex, its causal diamonds are the
compact diamonds of $M$, and $N$ is globally hyperbolic (MS Def.~3.70). In
$M'$,
$D_{\eta_\lambda}(\Sigma')\subseteq D_{g'}(\Sigma')\subseteq D_{\eta_\mu}(\Sigma')$
(narrower cones give larger domains of dependence). The first inclusion puts
every relative-interior point of $\Sigma'$, in particular $x$ and $b$, in the
interior; the second puts the projection of every interior point with $t\neq0$
into $\operatorname{int}B_0$, so its vertical segment to $\{t=0\}$ lies in
$\Ichr(q^+)\cap\Ifut(z^-)\subseteq N'$ for suitable $q^+$, $z^-$, joining it
inside $N'$ to the connected set $\operatorname{int}\Sigma'$. Transfer by (ii).
\end{proof}

By Lemma~\ref{lem:D}, $N$ is an object of $\GH$ and the inclusion $N\hookrightarrow
M$ and $\psi|_N:N\to M'$ are morphisms. By (LC),
\begin{equation}\label{eq:transfer}
x\,B_M\,b\iff x\,B_N\,b\iff\psi(x)\,B_{M'}\,\psi(b).
\end{equation}

\begin{lemma}[Flat detour]\label{lem:F}
Let $P,Q$ lie in the flat region of $M'$, be $\eta_\lambda$-spacelike, and be
joined by a smooth embedded arc $\kappa$ contained in a hyperplane
$H=\{t=f(s)\}$ with $f$ affine, $|f'|<1/\lambda$, and at positive distance from
$\overline{T_{3r}}$. Then $P,Q$ are spacelike-separated in $M'$, and there is a
compact acausal slab $\Sigma''\subset H$, containing $P,Q$ in its relative
interior, such that $N''=\operatorname{int}D_{M'}(\Sigma'')$ lies in the flat
region, equals $\operatorname{int}D_\Mink(\Sigma'')$ under the identity, and is
an object of $\GH$ whose inclusions into $M'$ and $\Mink$ are morphisms. Hence
$P\,B_{M'}\,Q\iff P\,B_\Mink\,Q$.
\end{lemma}

\begin{proof}
A $g'$-causal curve between $P,Q$ would be $\eta_\lambda$-causal. Along
$g'$-causal curves $\tfrac{d}{d\tau}(t-f(s))\ge\dot t(1-\lambda|f'|)>0$ and
$t-f(s)\to\pm\infty$, so $H$ is a Cauchy hypersurface for $g'$ and for $\eta$.
Extend $\kappa$ slightly past $P$ and $Q$ and let $\Sigma''$ be its closed
$\delta$-neighbourhood in $H$, with $\delta$ small compared with the distance from
$\kappa$ to $\overline{T_{3r}}$ and with the curvature radius of $\kappa$ in $H$.
A Lorentz boost (an isometry of $\eta$) makes $H$ horizontal; there the metric
is exactly $\eta$ on a $C\delta$-neighbourhood of $\Sigma''$, and the argument of
Lemma~\ref{lem:D}, with the straight tube replaced by the $2\delta$-tube about the
boosted arc and outward escape along the gradient of the distance to the arc,
applies with $(M,M')$ replaced by $(M',\Mink)$. It shows that both domains of
dependence lie in that neighbourhood, that they coincide, and that their
interiors are connected, causally convex and globally hyperbolic.
\end{proof}

\begin{proof}[Proof of Theorem \ref{thm:main}]
(a) Suppose $B_\Mink\neq\U$; by Theorem~\ref{thm:flat}, $B_\Mink\subseteq\Jm$
relates no spacelike pair. Let $x\,B_M\,b$ with $x,b$ spacelike-separated, and
transfer by \eqref{eq:transfer}. Take $w=(-4r,0,0,0)$, $v=(4r,L,0,0)$ and the arc
$t(s)=-4r+(8r/L)s$, $y(s)=(10r\sin(\pi s/L),0)$, $0\le s\le L$: where $|t|\le3r$,
$s/L\in[1/8,7/8]$ and $|y|\ge10r\sin(\pi/8)>3r$, elsewhere $|t|>3r$, and the
slope $8r/L<1/\lambda$. Vertical segments are $\eta_\mu$-timelike, so
$w\ll'\psi(x)$ and $\psi(b)\ll'v$; (CB) gives $w\,B'\,\psi(x)$ and
$\psi(b)\,B'\,v$, and (T) gives $w\,B'\,v$. Lemma~\ref{lem:F} gives $w\,B_\Mink\,v$
with $w,v$ spacelike, a contradiction. If $x\,B_M\,b$ with
$x\in\Jp(b)\setminus\{b\}$, Lemma~\ref{lem:R} gives $e\ll x$ spacelike to $b$, and
$e\,B_M\,b$ by (CB) and (T), reducing to the previous case.

(b) Suppose $B_\Mink=\U$. For spacelike $x,b\in M$, transfer to $M'$ and use
$x'=(4r,0,0,0)$, $b'=(-4r,L,0,0)$ with the arc $t(s)=4r-(8r/L)s$, same $y(s)$:
$\psi(x)\,B'\,x'$ and $b'\,B'\,\psi(b)$ by (CB), and $x'\,B'\,b'$ by
Lemma~\ref{lem:F}, since $B_\Mink=\U$; so $x\,B_M\,b$. By Lemma~\ref{lem:E} and
(T), every pair is related.
\end{proof}

\section{Causal sets: an exact Stein theorem}\label{sec:causets}

Let $\mathsf{CS}$ have locally finite posets as objects and order-embeddings
with order-convex image as morphisms. A becoming theory on $\mathsf{CS}$ is
defined as before, with (CB) read as $x\le b\Rightarrow x\,B\,b$.

\begin{proposition}\label{prop:causets}
Every becoming theory on $\mathsf{CS}$ is either the causal order in every
causal set or the universal relation in every causal set.
\end{proposition}

\begin{proof}
If $x,b$ are incomparable, $\{x,b\}$ is convex, so both embeddings of the
two-element antichain are morphisms: the verdict on incomparable pairs is a
single symmetric bit $\beta$, the same in every causal set. If $x>b$, the
interval $I=[b,x]$ is finite and convex, so $x\,B\,b$ depends only on $I$. If
$\beta=1$, in $I\sqcup\{z\}$ with $z$ incomparable to all, $x\,B\,z\,B\,b$; so
every such verdict is positive and $B$ is universal. If $\beta=0$ and
$x\,B\,b$ for some $x>b$, adjoin $e$ with $e<x$ only: $I$ remains convex, so
$x\,B\,b$ there, $e\,B\,x$ by (CB), and $e\,B\,b$ with $e$ incomparable to $b$,
contradicting $\beta=0$.
\end{proof}

Wüthrich and Callender state that ``The analogue of this theorem, however, is
straightforwardly false in CST'' \citep[\S3]{WC2017}. Their counterexample, on
the diamond causal set $d<b,c<a$, adds the incomparable pair $\{b,c\}$ to the
causal order and is invariant under that causal set's automorphisms. It is not
the restriction of any becoming theory in the present sense: relating $b$ to
$c$ fixes $\beta=1$, and Proposition~\ref{prop:causets} then relates $b$ to $d$,
which their relation excludes. Their own footnote~5 observes that the
counterexample ``could be a small proper part of a much larger causet''; (LC) is
exactly the requirement that the relation on such a part agree with the
relation in the larger causal set. The diagnosis matches the continuum case:
invariance within one structure is too weak, and locality across structures
restores Stein's theorem, here with no residual freedom because a causal set
has no light-cone boundary. Their claim is true on its own terms; the present
result uses a stronger, cross-structure hypothesis. Note also that
$\mathsf{CS}$ includes posets that are not manifoldlike (antichains, disjoint
unions), and the proof uses them: this is the causal-set counterpart of the
off-shell premise of Remark~\ref{rem:offshell}.

\section{Gödel universes are isolated}\label{sec:godel}

\begin{proposition}\label{prop:godel}
Extend $\GH$ to all oriented, time-oriented spacetimes with the same
morphisms. If $M$ is totally vicious ($p\ll p$ for all $p$), then every
morphism into $M$ is onto, and every theory satisfying (CB) has $B_M=\U_M$.
\end{proposition}

\begin{proof}
In a totally vicious spacetime $\Ichr(p)=\Ifut(p)=M$ for every $p$ (MS \S3.1).
Let $V$ be nonempty, open and causally convex, $p\in V$ and $q\in M$. There are
timelike curves $p\to q$ and $q\to p$; their concatenation is a causal curve
with both endpoints in $V$, so it lies in $V$, and $q\in V$. Thus $V=M$
(compare MS Remark~2.12).
\end{proof}

Gödel's universe is totally vicious, so universal becoming there is forced;
this is Gödel's first step, in Stein's terms. But local covariance transmits
constraints only along morphisms, and the only morphisms into a totally
vicious spacetime are its own isometries. Gödel's modal step, from what holds
in such a universe to what holds in ours \citep[pp.~561--562]{Godel1949}, as
quoted in \citep[\S6]{Monton2006}, requires that becoming be the same
feature across nomologically possible universes. Local covariance is one
precise way of requiring a rule to be the same across universes, and it does
not transmit anything out of totally vicious ones by morphisms into them. (We
have not shown that no morphisms lead out of them; Proposition~\ref{prop:godel}
concerns morphisms into $M$.) A version of Gödel's step that does go through
would need a different cross-universe link. Proposition~\ref{prop:nonlocal} shows what such links
look like: natural but non-local. Compare the conceptual responses of
\citet{Savitt2005}, \citet{Dieks2006} and \citet{Aames2022}, which locate
becoming locally along worldlines.

\section{Energy conditions}\label{sec:energy}

Let $\GH_{\mathrm{DEC}}$ be the full subcategory of spacetimes whose Einstein
tensor satisfies the dominant energy condition. It is closed under restriction
to causally convex open subsets and contains $\Mink$.

\begin{proposition}[The auxiliary spacetime is unavailable]\label{prop:dec}
If $M\in\GH_{\mathrm{DEC}}$ is flat on a neighbourhood of some Cauchy surface,
then $M$ is flat. In particular $M'$ above lies in $\GH_{\mathrm{DEC}}$ only if
$g$ is flat on $T_{2r}$.
\end{proposition}

\begin{proof}
On the Cauchy surface the metric is flat and the Einstein tensor vanishes. By the conservation theorem
for matter obeying the dominant energy condition \citep[\S4.3]{HE1973}, the
Einstein tensor vanishes on its domain of dependence, i.e.\ everywhere. For
$M'$, any slice $\{t=c\}$ with $c<-3r$ has a flat neighbourhood.
A vacuum globally hyperbolic spacetime with flat data $(\Sigma,h,0)$ on a
Cauchy surface is a development of those data and so embeds isometrically in
their maximal development \citep[Thm.~3]{CBG1969}. That maximal development is
the flat static spacetime $\mathbb R\times\Sigma$, $-dt^2+h$ (not Minkowski in
general; e.g.\ $\Sigma=T^3$): it is a development, and if a development
properly extended it, a point $p$ outside its image would lie on a timelike
geodesic from $\Sigma$ whose image under the extension coincides with the
corresponding complete geodesic of $\mathbb R\times\Sigma$, so $p$ would lie in
the image. Hence the spacetime is flat.
\end{proof}

The same reasoning constrains any repair that keeps the original spacetime and
inserts small flat pockets near $w$, $v$ and the detour: in a vacuum region
whose causal past on some Cauchy surface is vacuum and unmodified, the
conservation theorem forces the modified spacetime to be vacuum there with the
same data, so no pocket can be inserted. Pockets must be fed by matter brought
in from the past and routed around $D(\Sigma)$ afterwards. The proof of
Theorem~\ref{thm:main} goes through verbatim in $\GH_{\mathrm{DEC}}$ given
the following sufficient condition, which we leave open.

\begin{quote}
\textbf{Pocket problem.} Given $M\in\GH_{\mathrm{DEC}}$ and $\Sigma$ as above,
find $M''\in\GH_{\mathrm{DEC}}$ containing an isometric copy of a neighbourhood
of $D_M(\Sigma)$ whose image is a domain of dependence, and exactly flat
causally convex regions containing points $w\ll x$ and $v\gg b$ joined by the
flat domain of dependence of a slab.
\end{quote}

The constraint is an inequality, not a field equation, so constructions can be
kinematic, and explicit models with flat regions chronologically adjacent to
curved dominant-energy regions exist (Vaidya spacetimes before an incoming, or
after an outgoing, null-dust shell). Thin shells will not suffice in general,
since null junctions require matching intrinsic geometry. A solution would
give Theorem~\ref{thm:main} for spacetimes with physically reasonable matter;
a vacuum version would require bounded interior flat regions adjacent to
arbitrary curved vacuum regions, which we have not found in the gluing
literature (an abstract-level survey: Carlotto--Schoen, Corvino--Schoen,
Aretakis--Czimek--Rodnianski, Czimek--Rodnianski).

\section{What this shows about becoming}\label{sec:phil}

\paragraph{The robust part of Stein's theorem.} In general relativity, under
locality, nothing outside an event's causal past can have become for it
unless everything has. This holds in every globally hyperbolic spacetime,
however asymmetric, and does not rely on any foliation. Those who want an
extended, spacelike ``now'' as part of objective becoming face Stein's
dilemma in general relativity too, unless they give up locality or enrich
the spacetime structure itself, adding a preferred foliation or reference
field so that becoming is local relative to the richer structure (in effect
Monton's move). The theorem then does not apply, because the category, and
with it the notion of an isometric embedding, has changed; the cost is the
added structure.

\paragraph{The fragile part.} Uniqueness does not survive. Minkowskian
symmetry forced a choice between $\Ichr$ and $\Jm$; general relativity leaves
open how lightlike-separated events stand, and a rule may let curvature along
the light ray decide. Clifton and Hogarth's debate over causal versus
chronological becoming \citep[pp.~364--365]{CH1995} thus acquires a third
family of options; the choice between the endpoints is equivalent to a
regularity requirement (Proposition~\ref{prop:reg}).

\paragraph{Locality does the work symmetry did.} In a single curved spacetime,
definability-style invariance constrains almost nothing
(Proposition~\ref{prop:single}); natural global rules such as cosmological
time violate only locality (Proposition~\ref{prop:nonlocal}). So the
curved-spacetime case for ``only the past has become'' rests on a locality
principle of the kind locally covariant field theory takes for granted and
akin to Viglione's Causal Connection Condition \citep{Viglione2026} (which is
stricter: (LC) allows dependence on spacelike events inside the convex region). Rejecting
locality is a coherent option, but it is the option of global, cosmological
becoming, with its own costs \citep{Dieks2006,Viglione2026}.

\paragraph{Limits.} The theorem concerns relations of becoming, not
ontology; it does not decide between presentism, the growing block and
eternalism. It assumes that becoming is defined for every globally hyperbolic
geometry (Remark~\ref{rem:offshell}); Section~\ref{sec:energy} shows that
dropping this premise is not routine. Spacetimes that are not globally hyperbolic are not treated beyond
Proposition~\ref{prop:godel}.

\section{Prior work}\label{sec:prior}

Stein's theorem and its Clifton--Hogarth form are the direct predecessors
\citep{Stein1968,Stein1991,CH1995}. Point-based presents \citep{Rakic1997,Muller2006}
and the causal-diamond present of Arthur and Savitt
\citep{Arthur2006,Arthur2019,Savitt2009,Savitt2015}, with Dorato's critique
\citep{Dorato2011} and Aames's application to Gödel universes \citep{Aames2022},
are conceptual proposals without a curved-spacetime classification. In flat
spacetime, definability of sets from relativistic structure is governed by
automorphism groups \citep{MSS2025}. A citation sweep of
\citet{CH1995} and \citet{Stein1991} (Semantic Scholar, 24 October 2026), and
full-text reading of the works named here, found no prior extension of
Stein's theorem to curved spacetime. This bounds, but does not certify,
priority.

\paragraph{Note on method.} This manuscript was developed with substantial
assistance from large language models (Anthropic Claude and OpenAI Codex),
including literature search, proof drafting and an independent cross-model
proof check that found and repaired an error in an earlier draft. Quotations
were checked against the primary texts; responsibility for errors is the
author's.

\begin{thebibliography}{99}
\bibitem[Aames(2022)]{Aames2022} J. Aames, Temporal becoming in a relativistic
universe: causal diamonds and Gödel's philosophy of time, \emph{Eur.\ J.\
Philos.\ Sci.} 12 (2022) 44. doi:10.1007/s13194-022-00471-z.
\bibitem[Andersson et al.(1998)]{AGH1998} L. Andersson, G. J. Galloway and
R. Howard, The cosmological time function, \emph{Class.\ Quantum Grav.} 15
(1998) 309--322. doi:10.1088/0264-9381/15/2/006.
\bibitem[Arthur(2006)]{Arthur2006} R. T. W. Arthur, Minkowski spacetime and
the dimensions of the present, in D. Dieks (ed.), \emph{The Ontology of
Spacetime}, Elsevier, 2006, ch.~7. doi:10.1016/s1871-1774(06)01007-2.
\bibitem[Arthur(2019)]{Arthur2019} R. T. W. Arthur, Relativity and the
present, in \emph{The Reality of Time Flow}, Springer, 2019, ch.~6.
doi:10.1007/978-3-030-15948-1\_6.
\bibitem[Bernal and Sánchez(2006)]{BS2006} A. N. Bernal and M. Sánchez,
Further results on the smoothability of Cauchy hypersurfaces and Cauchy time
functions, arXiv:gr-qc/0512095.
\bibitem[Brunetti et al.(2003)]{BFV2003} R. Brunetti, K. Fredenhagen and
R. Verch, The generally covariant locality principle, \emph{Commun.\ Math.\
Phys.} 237 (2003) 31--68. doi:10.1007/s00220-003-0815-7.
\bibitem[Choquet-Bruhat and Geroch(1969)]{CBG1969} Y. Choquet-Bruhat and
R. Geroch, Global aspects of the Cauchy problem in general relativity,
\emph{Commun.\ Math.\ Phys.} 14 (1969) 329--335. doi:10.1007/bf01645389.
\bibitem[Clifton and Hogarth(1995)]{CH1995} R. Clifton and M. Hogarth, The
definability of objective becoming in Minkowski spacetime, \emph{Synthese} 103
(1995) 355--387. doi:10.1007/BF01089733.
\bibitem[Dieks(2006)]{Dieks2006} D. Dieks, Becoming, relativity and locality,
in D. Dieks (ed.), \emph{The Ontology of Spacetime}, Elsevier, 2006, ch.~8.
doi:10.1016/s1871-1774(06)01008-4.
\bibitem[Dorato(2011)]{Dorato2011} M. Dorato, The Alexandroff present and
Minkowski spacetime: why it cannot do what it has been asked to do, in
\emph{Explanation, Prediction, and Confirmation}, Springer, 2011.
doi:10.1007/978-94-007-1180-8\_26.
\bibitem[Fewster and Verch(2012)]{FV2012} C. J. Fewster and R. Verch,
Dynamical locality and covariance: what makes a physical theory the same in
all spacetimes?, arXiv:1106.4785.
\bibitem[Gödel(1949)]{Godel1949} K. Gödel, A remark about the relationship
between relativity theory and idealistic philosophy, in P. A. Schilpp (ed.),
\emph{Albert Einstein: Philosopher-Scientist}, 1949, pp.~557--562.
\bibitem[Hawking and Ellis(1973)]{HE1973} S. W. Hawking and G. F. R. Ellis,
\emph{The Large Scale Structure of Space-Time}, Cambridge University Press,
1973. doi:10.1017/cbo9780511524646.
\bibitem[Madarász et al.(2025)]{MSS2025} J. Madarász, M. Stannett and
G. Székely, Definable coordinate geometries over fields, part 1: theory,
arXiv:2507.10279.
\bibitem[Minguzzi and Sánchez(2008)]{MS2008} E. Minguzzi and M. Sánchez, The
causal hierarchy of spacetimes, in \emph{Recent Developments in
Pseudo-Riemannian Geometry}, ESI Lectures in Mathematics and Physics, EMS,
2008, pp.~299--358. doi:10.4171/051-1/9;
arXiv:gr-qc/0609119.
\bibitem[Monton(2006)]{Monton2006} B. Monton, Presentism and quantum gravity,
in D. Dieks (ed.), \emph{The Ontology of Spacetime}, Elsevier, 2006, ch.~14
(read in the author's 2005 preprint).
\bibitem[Müller(2006)]{Muller2006} T. Müller, On the problem of defining the
present in special relativity: a challenge for tense logic, in
\emph{Time and History}, De Gruyter, 2006. doi:10.1515/9783110333213.441.
\bibitem[Rakić(1997)]{Rakic1997} N. Rakić, Past, present, future, and special
relativity, \emph{Brit.\ J.\ Philos.\ Sci.} 48 (1997) 257--280.
doi:10.1093/bjps/48.2.257.
\bibitem[Savitt(2005)]{Savitt2005} S. Savitt, Time travel and becoming,
\emph{The Monist} 88 (2005) 413--422. doi:10.5840/monist200588323.
\bibitem[Savitt(2009)]{Savitt2009} S. Savitt, The transient nows, in
\emph{Quantum Reality, Relativistic Causality, and Closing the Epistemic
Circle}, Springer, 2009. doi:10.1007/978-1-4020-9107-0\_18.
\bibitem[Savitt(2015)]{Savitt2015} S. Savitt, I $\heartsuit$ $\Diamond$s,
\emph{Stud.\ Hist.\ Philos.\ Mod.\ Phys.} 50 (2015) 19--24.
doi:10.1016/j.shpsb.2015.02.001.
\bibitem[Stein(1968)]{Stein1968} H. Stein, On Einstein--Minkowski space--time,
\emph{J.\ Philos.} 65 (1968) 5--23. doi:10.2307/2024512.
\bibitem[Stein(1991)]{Stein1991} H. Stein, On relativity theory and openness
of the future, \emph{Philos.\ Sci.} 58 (1991) 147--167. doi:10.1086/289609.
\bibitem[Viglione(2026)]{Viglione2026} F. Viglione, Not-so-absolute cosmic
simultaneity, \emph{Erkenntnis} (2026). doi:10.1007/s10670-026-01075-2.
\bibitem[Wüthrich and Callender(2017)]{WC2017} C. Wüthrich and C. Callender,
What becomes of a causal set?, \emph{Brit.\ J.\ Philos.\ Sci.} 68 (2017)
907--925. doi:10.1093/bjps/axv040; arXiv:1502.00018.
\end{thebibliography}
\end{document}
